% ============================================================
%  TurkLang 2026 --- Speech notes for the 14-minute talk
%  Paper: Neural Morphosyntactic Tagging and Dependency Parsing for Uzbek
%  Author: Sanatbek Matlatipov
%
%  Optimised for printing on a black & white printer:
%    - 12pt serif body, 1.4 line spacing
%    - bold cues only (no colours), clear slide separators
%    - generous left margin for handwritten annotations
%    - cumulative timing budget printed next to each slide
%
%  Build with:  pdflatex speech_notes.tex
% ============================================================
\documentclass[12pt,a4paper]{article}

\usepackage[utf8]{inputenc}
\usepackage[T1]{fontenc}
\usepackage{lmodern}
\usepackage[margin=2.2cm,left=3.2cm,right=2.2cm]{geometry}
\usepackage{setspace}
\usepackage{enumitem}
\usepackage{titlesec}
\usepackage{fancyhdr}
\usepackage{microtype}
\usepackage{parskip}
\usepackage{amsmath,amssymb}

\setstretch{1.40}
\setlist[itemize]{leftmargin=1.3em,topsep=2pt,itemsep=2pt}

\titleformat{\section}
   {\normalfont\Large\bfseries\filright}
   {Slide \thesection.}{0.6em}{}
\titlespacing*{\section}{0pt}{1.2em}{0.4em}

\pagestyle{fancy}
\fancyhf{}
\fancyhead[L]{\small TurkLang 2026 --- Speech notes (14 min)}
\fancyhead[R]{\small S.~Matlatipov}
\fancyfoot[C]{\thepage}
\renewcommand{\headrulewidth}{0.4pt}

\newcommand{\timing}[2]{%
   \par\noindent\textbf{\small [#1\,$\to$\,#2\,min]}\par\vspace{0.2em}}
\newcommand{\say}[1]{\textbf{#1}}
\newcommand{\pause}{\par\vspace{0.4em}\textbf{[pause]}\par}

\begin{document}

% ------------------------------------------------------------
%  COVER PAGE
% ------------------------------------------------------------
\begin{center}
   {\LARGE\bfseries Speech notes --- 14 minutes}\\[0.6em]
   {\large Neural Morphosyntactic Tagging\\and Dependency Parsing for Uzbek}\\[0.4em]
   {\normalsize XIV International Conference on Computer Processing of Turkic Languages}\\
   {\normalsize \textbf{TurkLang 2026}}\\[1.2em]
   \textit{Sanatbek Matlatipov}\\
   National University of Uzbekistan, Tashkent
\end{center}

\vspace{1em}
\noindent\textbf{How to use these notes.}
Each section maps to one slide. The bracketed times
\textbf{[start\,$\to$\,end]} are cumulative running time. Phrases in
\textbf{bold} are the ones to emphasise; \textbf{[pause]} marks a short
breath. Plain text is the spoken script --- read naturally, not
word-for-word.

\vspace{0.5em}
\noindent\textbf{Total budget:} 14 minutes of speech + $\approx$1 minute
buffer for transitions. Target around 0:50--1:10 per slide.

\vfill\hrule\vspace{0.5em}
\noindent\textit{Print double-sided on A4. The wide left margin is for handwritten cues.}

\newpage

% ------------------------------------------------------------
%  SLIDE 1 --- Title
% ------------------------------------------------------------
\section*{Slide 1 --- Title}
\timing{0:00}{0:30}

Good morning / afternoon, and thank you to the TurkLang organisers for
the opportunity to present.

My name is \say{Sanatbek Matlatipov}, from the National University of
Uzbekistan. The talk is on \say{Neural Morphosyntactic Tagging and
Dependency Parsing for Uzbek} --- a companion piece to the gold-standard
treebank work, but here the focus is on \say{the model and the
evaluation}.

% ------------------------------------------------------------
%  SLIDE 2 --- Outline
% ------------------------------------------------------------
\section*{Slide 2 --- Outline}
\timing{0:30}{1:00}

The talk is structured in three blocks:

\begin{itemize}
   \item First, \textbf{why} Uzbek is hard, and what existing work has done.
   \item Second, the \textbf{pipeline} --- preprocessing, representation, and the joint tagging / parsing model.
   \item Third, the \textbf{results}, the error analysis, and what they tell us about Uzbek parsing.
\end{itemize}

% ------------------------------------------------------------
%  SLIDE 3 --- Motivation
% ------------------------------------------------------------
\section*{Slide 3 --- Why Uzbek morphosyntactic analysis is hard}
\timing{1:00}{2:00}

Uzbek is a \say{Turkic, agglutinative, free-SOV language} with roughly
35 million speakers. Three properties make UD analysis substantially
harder than for analytic languages.

\begin{itemize}
   \item \textbf{Long suffix chains} on both verbs and nouns.
   \item Rich \textbf{case, possessive, and tense-aspect-mood} feature bundles.
   \item Which together produce \textbf{enormous lexical sparsity} --- many surface forms, but each appears very few times.
\end{itemize}

\noindent
And the UD treebank we evaluate on, \say{UzUDT}, has \textbf{under ten
thousand tokens}. So our research question is: \say{can a reproducible
neural pipeline --- Apertium normalization plus BERTbek embeddings ---
deliver a strong baseline despite all of this?}

% ------------------------------------------------------------
%  SLIDE 4 --- Related work
% ------------------------------------------------------------
\section*{Slide 4 --- Where this work fits}
\timing{2:00}{2:45}

Briefly on related work. The framework is \textbf{Universal Dependencies}.
For neural pipelines I build on \textbf{Stanza}, the \textbf{biaffine
parser of Dozat and Manning}, and \textbf{UDify}. For Turkic in general
there is excellent work on Turkish, Kazakh, Uyghur and Tatar; for Uzbek
specifically there is the original \say{Uzbek-UT} from
Akhundjanova and Talamo, and the newer \say{UzUDT} that I use as my
benchmark.

\pause

Two further ingredients are the \textbf{Apertium} finite-state analyser
for Uzbek --- through the work of Salaev, Sharipov and others --- and the
\textbf{BERTbek} monolingual Transformer by Kuriyozov. Combining these
two has not been done in a reproducible UD-aligned baseline before, and
that is what this paper contributes.

% ------------------------------------------------------------
%  SLIDE 5 --- Data
% ------------------------------------------------------------
\section*{Slide 5 --- UzUDT: data and splits}
\timing{2:45}{3:30}

The data is the publicly available \say{UzUDT} treebank --- 681 sentences
and 7\,542 tokens, annotated in UD version 2.

\begin{itemize}
   \item I take the original test set as is: \textbf{198 sentences, 2\,101 tokens}.
   \item And I deterministically carve a 10\% \textbf{development set} from the training portion --- 48 sentences, 518 tokens --- so I can do early stopping and model selection without touching the test set.
\end{itemize}

\noindent
This leaves \textbf{435 training sentences and roughly 4\,900 training
tokens}. It is \say{a small treebank by any measure} --- which is part of
why the choice of representation matters.

% ------------------------------------------------------------
%  SLIDE 6 --- Pipeline overview
% ------------------------------------------------------------
\section*{Slide 6 --- End-to-end pipeline}
\timing{3:30}{4:30}

This slide shows the full architecture. From left to right, four stages:

\begin{itemize}
   \item \textbf{Apertium} morphology-aware normalization on raw text.
   \item \textbf{BERTbek} encoding on the resulting subword sequence.
   \item \textbf{Super-token fusion} that aggregates subwords back to UD word boundaries.
   \item And a \textbf{joint head} that does tagging --- UPOS, XPOS, UFeats --- and biaffine dependency parsing at the same time.
\end{itemize}

\noindent
The key design point is that we \say{never break the UD word
alignment}: every component upstream of the parser keeps a stable
subword-to-word mapping. That makes evaluation honest, and it makes the
model reproducible across runs.


% ------------------------------------------------------------
%  SLIDE 7 --- Preprocessing
% ------------------------------------------------------------
\section*{Slide 7 --- Morphology-aware preprocessing}
\timing{4:30}{5:30}

Preprocessing has three steps:

\begin{itemize}
   \item \textbf{Text normalization} --- Unicode and UD-compliant tokenisation, so token boundaries are stable.
   \item \textbf{Apertium morphological analysis} --- lexicon lookup and morphotactic parsing. This \say{stabilises lemmas across inflectional variants} and dramatically reduces lexical sparsity.
   \item \textbf{Subword tokenisation} for BERTbek, while preserving the UD word boundaries.
\end{itemize}

\noindent
The figure on the right shows what \say{Uzbek morphotactics} actually
looks like --- long suffix chains, sometimes five or six morphemes on a
single verb or noun. Apertium handles these chains explicitly, so the
neural model does not have to memorise every surface form from a few
hundred sentences.

% ------------------------------------------------------------
%  SLIDE 8 --- Token representations
% ------------------------------------------------------------
\section*{Slide 8 --- Token representations and super-token fusion}
\timing{5:30}{7:00}

BERTbek produces \textbf{one vector per WordPiece}, but UD evaluation
operates at \textbf{word level}. We bridge that gap with \say{super-token
fusion} --- a fancy name for a simple operation. For each UD word $w_i$,
we take its $m_i$ subword vectors and \textbf{mean-pool} them into a
single representation $\mathbf{e}_{\mathrm{ctx}}(w_i)$.

\pause

The final per-token vector concatenates \textbf{five channels}: a word
embedding, a character embedding, a UPOS embedding using the
\emph{predicted} (not gold) tag --- so no oracle leakage --- an optional
explicit morphological channel which is \textbf{disabled in this
baseline}, and finally the BERTbek mean-pooled context vector.

\pause

Why mean pooling rather than something fancier? Two reasons. First, it
is \textbf{deterministic and language-agnostic}; nothing to tune.
Second, in our experiments it was \textbf{competitive with last-subword
pooling} for tagging, while avoiding a positional heuristic. We want a
clean baseline that other groups can reproduce.

% ------------------------------------------------------------
%  SLIDE 9 --- Joint tagging and parsing
% ------------------------------------------------------------
\section*{Slide 9 --- Joint tagging and biaffine parsing}
\timing{7:00}{8:00}

A \textbf{BiLSTM} consumes the per-token vectors and produces contextual
hidden states. From those, three softmax heads predict
\textbf{UPOS, XPOS, and UFeats}, and a \textbf{biaffine parser} --- the
classic Dozat--Manning architecture --- predicts arcs and relation labels.

\pause

The training objective is the boxed equation: \say{the tagging loss plus
$\lambda$ times the parsing loss}, where the parsing loss is itself the
sum of arc and relation cross-entropies. We set $\lambda = 1.0$ --- a
deliberate choice to keep the baseline simple. Optimisation is AdamW
with a peak learning rate of \textbf{two times ten to the minus five},
\textbf{linear warmup over the first ten percent of steps}, batch size
32, and early stopping on the development set.

% ------------------------------------------------------------
%  SLIDE 10 --- Hyperparameters
% ------------------------------------------------------------
\section*{Slide 10 --- Hyperparameters at a glance}
\timing{8:00}{8:45}

I will not read this table line by line --- it is in the paper. The
takeaway is that \textbf{nothing was hand-tuned per run}. Every
hyperparameter is fixed and reported. That is the price you pay to
claim \say{a reproducible baseline}, and it is a price worth paying for
a low-resource language where comparing future improvements honestly
matters more than chasing a fragile peak.

% ------------------------------------------------------------
%  SLIDE 11 --- Tagging results
% ------------------------------------------------------------
\section*{Slide 11 --- Morphosyntactic tagging results}
\timing{8:45}{9:45}

The tagging results, read carefully.

\begin{itemize}
   \item \textbf{Coarse tags are robust.} UPOS reaches \textbf{86.1\%} and XPOS \textbf{84.0\%}.
   \item \textbf{Fine-grained features are harder.} UFeats drops to \textbf{70.1\%}, and AllTags --- the metric that demands every feature correct for a token to count --- drops to \textbf{61.4\%}.
\end{itemize}

\noindent
The drop is \say{driven by suffix-bundle sparsity}. A finite verb in
Uzbek can carry tense, aspect, mood, person and number all at once. If
even one of those is wrong, AllTags marks the token as a miss. There
are also some residual annotation inconsistencies in the underlying
treebank --- small in number, but they hit the strictest metric hardest.

% ------------------------------------------------------------
%  SLIDE 12 --- Parsing results
% ------------------------------------------------------------
\section*{Slide 12 --- Dependency parsing results}
\timing{9:45}{10:45}

The parsing results, in the standard CoNLL-U metrics.

\begin{itemize}
   \item \textbf{UAS 74.2 / LAS 66.9} --- solid attachment, weaker relation labelling.
   \item The \textbf{seven-point UAS--LAS gap} tells a consistent story: \say{the parser knows where words attach, but it sometimes labels the edge incorrectly}.
   \item \textbf{CLAS, MLAS, BLEX} all fall in the 44--47 range, which is expected --- these stricter metrics multiply tagging errors and parsing errors together. They confirm that the difficulty lives \textbf{at the morphology--syntax interface}, not in either layer alone.
\end{itemize}

% ------------------------------------------------------------
%  SLIDE 13 --- Relation distribution
% ------------------------------------------------------------
\section*{Slide 13 --- What the parser has to predict}
\timing{10:45}{11:30}

This bar chart shows the \textbf{top 20 dependency relations} in UzUDT.
The distribution is \say{heavily skewed} --- punctuation, oblique,
nominal subject, and root absorb most of the mass. Mid-frequency
relations like \texttt{advcl}, \texttt{acl}, and \texttt{xcomp} carry
\textbf{the harder syntactic phenomena}.

\pause

The implication is direct: the dominant classes drive UAS, but they do
not push LAS up. \textbf{LAS is bounded above by how well the model
handles the long tail}, and that tail is data-starved.

% ------------------------------------------------------------
%  SLIDE 14 --- Discussion / error analysis
% ------------------------------------------------------------
\section*{Slide 14 --- Why Uzbek remains challenging}
\timing{11:30}{12:30}

Three factors persist beyond what subword context can repair.

\begin{itemize}
   \item \textbf{Subword boundaries do not align with morphemes.} WordPiece is trained on raw text, not Uzbek morphology, and grammatical information ends up scattered across subwords inconsistently.
   \item \textbf{Polyfunctional markers.} The same surface suffix can mark \texttt{case}, or play a \texttt{mark} role, or even encode discourse --- and surface form alone is not enough to disambiguate.
   \item \textbf{Complex predicates.} Light verbs and serial-verb constructions create representational ambiguity between \texttt{xcomp}, \texttt{ccomp}, and \texttt{compound:lvc}.
\end{itemize}

\pause

The errors concentrate on \textbf{obl versus obj} for case-marked
nominals, and on \textbf{advcl versus xcomp} for embedded clauses.
Apertium normalisation and super-token fusion \textbf{do not solve
these} --- they reduce the noise floor and make the baseline reproducible.
Solving them requires either \say{more annotated data} or \say{a richer
morpheme-aware tokenizer}.

% ------------------------------------------------------------
%  SLIDE 15 --- Conclusion
% ------------------------------------------------------------
\section*{Slide 15 --- Conclusion}
\timing{12:30}{13:30}

To summarise:

\begin{itemize}
   \item We deliver a \textbf{reproducible neural baseline} for Uzbek UD analysis on UzUDT, combining a Stanza-style biaffine model with Apertium normalisation and BERTbek super-token fusion.
   \item Tagging: \textbf{UPOS 86.1, XPOS 84.0, UFeats 70.1, AllTags 61.4}.
   \item Parsing: \textbf{UAS 74.2, LAS 66.9}.
   \item Errors are concentrated \textbf{at the morphology--syntax interface}, exactly where small treebanks are weakest.
\end{itemize}

\noindent
The natural next steps are: turn on the \textbf{explicit Apertium
feature channel}, scale to a \textbf{multi-domain Uzbek UD treebank},
add \textbf{cross-lingual transfer} from Turkish or Kazakh, and
experiment with \textbf{morpheme-aware tokenizers} for transformer
parsers like XLM-R or TURNA.

% ------------------------------------------------------------
%  SLIDE 16 --- Thank you
% ------------------------------------------------------------
\section*{Slide 16 --- Thank you}
\timing{13:30}{14:00}

That brings me to the end. The code, configurations, and trained models
are available on request. I would like to thank the open-source teams
behind \textbf{Stanza, Apertium, and BERTbek} --- without them this kind
of reproducible baseline for a low-resource language would simply not
be possible.

\pause

Thank you very much --- I am happy to take your questions.

% ------------------------------------------------------------
%  Q & A prep --- printed as a back page for the speaker only
% ------------------------------------------------------------
\newpage
\section*{Likely Q\&A --- quick cues}

\noindent\textbf{Q. Why disable the explicit morphological feature channel?}\\
We wanted a \textbf{clean baseline}. With it disabled, every gain from
future morphology-aware variants is unambiguously attributable. Turning
it on is listed as immediate future work.

\vspace{0.4em}
\noindent\textbf{Q. Why mean pooling rather than attention or last subword?}\\
\textbf{Reproducibility and simplicity}. Mean pooling has \textbf{no
trainable parameters}, no positional heuristic, and was empirically
competitive with last-subword in our development experiments.

\vspace{0.4em}
\noindent\textbf{Q. Why not a full transformer parser end-to-end?}\\
The Stanza-style pipeline is \textbf{lightweight, fast on CPU, and
trivially comparable} with existing UD baselines. A full transformer
parser is excellent future work but a different paper.

\vspace{0.4em}
\noindent\textbf{Q. Could mBERT or XLM-R outperform BERTbek here?}\\
On a \textbf{very small treebank} like UzUDT, monolingual BERTbek tends
to beat multilingual baselines on downstream Uzbek tasks. With more
data and joint Turkic training, multilingual models become competitive.

\vspace{0.4em}
\noindent\textbf{Q. How much variance is in UAS / LAS across seeds?}\\
The paper reports a \textbf{single-seed} result. Multi-seed variance was
not the focus of this work, but it is a fair point: a small treebank
like UzUDT does carry non-trivial seed sensitivity, and a multi-seed
study is part of the planned follow-up.

\vspace{0.4em}
\noindent\textbf{Q. Is the 86.10 UPOS comparable to Turkish baselines?}\\
\textbf{Caution: different treebanks, different domain.} The number is
in the same range as small-Turkic UD baselines without contextual
embeddings, and \textbf{lags behind larger-treebank Turkish UPOS} (which
exceeds 92\%). The gap reflects training-set size much more than model
architecture.

\vspace{0.4em}
\noindent\textbf{Q. Does Apertium actually help, or is it noise?}\\
The paper motivates Apertium normalisation as a way to \textbf{stabilise
lemmas and reduce lexical sparsity} for an agglutinative language; the
reported results use that pipeline throughout. A \textbf{controlled
ablation} with/without Apertium is not reported in this paper and is
listed as future work.

\vspace{0.4em}
\noindent\textbf{Q. Will the code and configs be released?}\\
Yes --- \textbf{available on request} and being prepared for public
release alongside the UzUDT companion paper.

\end{document}
